\documentclass[11pt,a4paper]{article}
\usepackage[margin=19mm,headheight=15pt]{geometry}
\usepackage{fontspec}
\setmainfont{Noto Serif}
\usepackage{amsmath,amssymb,mathtools}
\usepackage{xcolor}
\usepackage{fancyhdr}
\usepackage{enumitem}
\usepackage{microtype}
\setlength{\parindent}{0pt}
\setlength{\parskip}{6pt}
\definecolor{navy}{HTML}{183153}
\definecolor{lightgray}{HTML}{F3F5F7}
\pagestyle{fancy}\fancyhf{}
\lhead{\small ME 5309 -- Mid-Semester Examination}
\rhead{\small ME 5309}
\cfoot{\thepage}
\begin{document}
\begin{center}
{\Large\bfseries ME 5309 -- Mid-Semester Examination}\\[3pt]
{\large Answer Booklet}\\[8pt]
\end{center}
\textbf{Instructions:} Attempt any three questions. Each question carries 10 marks. Show working where appropriate.
\section*{Question 1 \hfill [10 marks]}
\subsection*{Q1(a) \hfill [4 marks]}

\[\begin{aligned}f_x = 3200 pixels; field of view = 160 mm; 1 pixel = 0.125 mm\end{aligned}\]

\subsection*{Q1(b) \hfill [3 marks]}

The smallest resolvable force is 0.122 N.

\subsection*{Q1(c) \hfill [3 marks]}

\[\begin{aligned}F = 0.8187, G = 0.0906\end{aligned}\]

\section*{Question 3 \hfill [10 marks]}
\subsection*{Q3(a) \hfill [2 marks]}

\vspace{2mm}

\subsection*{Q3(b) \hfill [2 marks]}

The Tustin coefficients are y[k]=0.881765y[k-1]+0.0591175(x[k]+x[k-1]).
The DC gain is 0.0591175/(1-0.881765)=0.5.

\subsection*{Q3(c) \hfill [4 marks]}

\[\begin{aligned}L = [18,64]^T\end{aligned}\]

\subsection*{Q3(d) \hfill [2 marks]}

An observer uses the plant model and known control input to estimate unmeasured states such as velocity, and its poles can be designed.
A simple low-pass filter cannot exploit the state-space model in this way.
A filter may be safer when the model is poorly known or unreliable, because it relies less on model accuracy.\par\smallskip The observer uses the plant model and the measured output to reconstruct the unmeasured state. Its convergence depends on the model and the observer poles. A low-pass filter is simpler, but generally introduces phase lag and does not explicitly use the plant dynamics.

\section*{Question 4 \hfill [10 marks]}
\subsection*{Q4(a) \hfill [2 marks]}

Closed-loop characteristic polynomial is s²+(a+bk₂)s+(c+bk₁)=s²+4s+8, giving poles -2±2j.
At DC, θ/r = bNff/(c+bk₁), so unity gain requires Nff=(c+bk₁)/b=(6+2)/1=8.\par\smallskip Comparing coefficients with $s^2+4s+8$ gives $a+bk_2=4$ and $c+bk_1=8$. The feedforward gain is chosen from the nominal plant so the reference scaling is unity; here this gives $N_{ff}=8$. The closed-loop roots of $s^2+4s+8$ are $-2\pm2j$.

\subsection*{Q4(b) \hfill [2 marks]}

\[\begin{aligned}\theta_{ss} = 0.75\times8/(6+2)=0.75,\\ \text{the percentage tracking error is 25\%}\end{aligned}\]

\subsection*{Q4(c) \hfill [4 marks]}

With ξ̇=r-θ and u=-k₁θ-k₂θ̇+K\_Iξ, the homogeneous closed-loop characteristic polynomial is s³+(a+bk₂)s²+(c+bk₁)s+bK\_I.
Desired polynomial: (s+4)[(s+2)²+4]=(s+4)(s²+4s+8)=s³+8s²+24s+32.
For a=2,b=1,c=6, matching coefficients gives k₂=6, k₁=18 and K\_I=32.\par\smallskip First multiply $(s+4)(s^2+4s+8)=s^3+8s^2+24s+32$. Comparing this with the augmented characteristic polynomial gives $k_2=6$, $k_1=18$, and $K_I=32$. I checked the coefficients in descending powers of $s$.

\subsection*{Q4(d) \hfill [2 marks]}

Integral action removes steady-state error, although this only holds for the nominal values of b and c.
Saturation can cause windup, which should be handled by anti-windup.

\end{document}